% Context from chapters-tex/approximation.tex; for mathematical reference, not compilation.
\section{Approximation}

                                           
\subsection{Approximation in an Orthonormal Basis}

Let $\Bb = \{ \psi_m \}_m$ be an orthonormal basis of $\Ldeux([0,1]^d)$, with $d=1$ for signals or $d=2$ for images.
The complete expansion in this basis,
\eq{
	f = \sum_{m \in \ZZ} \dotp{f}{\psi_m} \psi_m
}
reconstructs the signal exactly. Approximation and processing methods modify the coefficients $\dotp{f}{\psi_m}$, introducing a reconstruction error.

The simplest approximation reconstructs the signal from a subset $I_M \subset \ZZ$ of $M$ coefficients:
\eq{
	f_M \eqdef \sum_{m \in I_M} \dotp{f}{\psi_m} \psi_m, 
	\qwhereq M = |I_M|.
}
The reconstructed signal $f_M$ is the orthogonal projection of $f$ onto the space
\eq{
	V_M \eqdef \Span\enscond{\psi_m}{m \in I_M}.
}
If the selected space $V_M$ depends on $f$, the overall approximation map $f \mapsto f_M$ may be nonlinear, even though projection onto each fixed space is linear.

Since the basis is orthogonal, the approximation error is 
\eq{
	\norm{f-f_M}^2 = \sum_{m \notin I_M} |\dotp{f}{\psi_m}|^2.
}
The central choice is the retained index set $I_M$, which may depend on the signal $f$.

                                           
\subsection{Linear Approximation}

Fixing $I_M$ independently of the input gives a linear approximation: the same coefficients are retained for every $f$. The map $f \mapsto f_M$ is then orthogonal projection onto a fixed space $V_M$, and it satisfies
\eq{
	(f+g)_M = f_M + g_M
}
For the Fourier basis and an even coefficient budget, one usually selects the low-frequency atoms
\eq{
	I_M = \{-M/2+1,\ldots,M/2\}.
}
For a 1-D wavelet basis, one usually selects the coarse wavelets
\eq{
	I_M=\{\text{coarsest scaling atoms and wavelets }(j,n)\text{ with }j\geq j_0\}
}
where $j_0$ is selected such that $|I_M|=M$.

\myfigure{
\tabtrois{
\image{approximation}{.32}{approximation-original}&
\image{approximation}{.32}{approximation-linear}&
\image{approximation}{.32}{approximation-non-linear}\\
Original $f$ & Linear $f_M^\ell$ & Nonlinear $f_M^n$
}
}{ 
	Linear versus nonlinear wavelet approximation.  
}{fig-approximation}

Figure \ref{fig-approximation}, center, illustrates linear wavelet approximation. Discarding all fine-scale details blurs singularities, where the approximation is least accurate.

                                           
\subsection{Nonlinear Approximation}

Allowing $I_M$ to depend on $f$ gives a nonlinear approximation. To choose $I_M$ that minimizes $\norm{f-f_M}$, orthogonality reduces the problem to retaining the $M$ coefficients of largest magnitude:
\eq{
	I_M= \{ M \text{ largest coefficients } |\dotp{f}{\psi_m}| \}.
}
When there are no ties at the cutoff, this can be obtained by thresholding
\eq{
	I_M = \enscond{m}{ |\dotp{f}{\psi_m}|>T }
}
where $T$ depends on the number of coefficients $M$, 
\eq{
	M = \#\enscond{m}{|\dotp{f}{\psi_m}|>T}.
}

  
\paragraph{Computation of the threshold.}

Order the coefficient magnitudes in nonincreasing order. For a finite vector, write
\eql{\label{eq-dfn-ordered-coefs}
	d_1\geq\cdots\geq d_N\geq0,\qquad\{d_m\}_{m=1}^{N}=\{|\dotp{f}{\psi_m}|\}_{m=1}^{N}
	.
}
If $d_M>d_{M+1}$, any threshold $d_{M+1}\leq T<d_M$ retains exactly $M$ coefficients. At a tie, a deterministic rule must select among equal coefficients. Thus the threshold $T$ and coefficient count $M$ do not determine each other uniquely. Figure~\ref{fig-decay-coefs} illustrates the relationship.

The following proposition shows that the decay of the ordered coefficients governs the decay of the nonlinear approximation error.

\begin{prop}\label{prop-equiv-comp-decay}
For $\alpha>0$, with $(d_m)_{m\geq1}$ in nonincreasing order, one has
\eql{\label{eq-decay-coef-approx}
	d_m = O( m^{-\frac{\al+1}{2}} )
	\quad\Longleftrightarrow\quad
	\norm{f-f_M}^2 = O(M^{-\al}).
}
\end{prop}

\begin{proof}
One has
\eq{
	\norm{f-f_M}^2 = \sum_{m>M} d_m^2
}
The forward implication follows from $\sum_{m>M}m^{-\alpha-1}\leq\alpha^{-1}M^{-\alpha}$. For the converse, taking $M$ even (integer rounding changes only constants),
\eq{
	d_M^2 \leq \frac{2}{M} \sum_{m=M/2+1}^{M} d_m^2 \leq 